\scanpage{51}
\uline{1} Automorphismes et repères des polygônes combinatoires. Polygône
combinatoire-type
\[
\begin{aligned}
\unsure{D_n|\mathbb D_n}&=\{\sigma,u\mid\sigma^2=\unsure{1\mathrm d|1},\sigma u\sigma^{-1}=u^{-1},u^n=1\}=\{\sigma,\sigma'\mid\sigma^2=\sigma'^{2}=1,(\sigma'\sigma)^n=1\}\\
D_\infty&=\{\sigma,u\mid\sigma^2=1,\sigma u\sigma^{-1}=u^{-1}\}=\{\sigma,\sigma'\mid\sigma^2=\sigma'^{2}=1\}=\mathbb Z_2*\mathbb Z_2
\end{aligned}
\]
2. Espaces affines -- Groupe affine $\mathrm{Aff}(E)$ ($\simeq\mathrm{Aff}(n,k)\simeq
\mathrm{Gl}(n).k^n$)

3. Métrique sur espaces affines et formes quadratiques.
Formes quadratiques non dégénérées.

Rappel des points développés la fois précédente :

a) Polygône géométrique comme couple $(S,A)$
\[
\begin{aligned}
S&\subset E\quad\text{(plan affine)}\\
A&\subset\mathrm{Dr}(E)
\end{aligned}
\]
Axiome : pour la relation d'incidence habituelle
$(s\in a)$, $(S,A)$ est un polygône combinatoire.

b) Réalisation géométrique d'un polygône
combinatoire $\Pi=(S,A,R)$
\diagram{\begin{tabular}{@{}c|l@{}}
\begin{tikzcd}[ampersand replacement=\&,row sep=small,column sep=large]
S\arrow[r,"\varphi_S"]\& E\\
A\arrow[r,"\varphi_A"]\& \mathrm{Dr}(E)
\end{tikzcd}&\shortstack[l]{compatibles avec\\\unsure{relation|relations} d'incidence}
\end{tabular}}
Si $\varphi_S(\unsure{s'|s})\ne\varphi_S(s)$ quand $s,s'\in S$ sont adjacents, \uline{alors}
$\varphi_A$ est connu quand on connaît $\varphi_S$ qui peut être
pris arbitrairement par ailleurs. On dira que
$\varphi$ est \deleted{\uline{\unsure{non dégénéré|non dégénérée}}} \inserted{\uline{fidèle}} si
\[
\varphi_S(s)\text{ incident à }\varphi_A(a)\Longrightarrow s\text{ incident à }a
\]
Cela implique $\varphi_S,\varphi_A$ injectifs et $\overline S=\varphi_S(S)$, $\overline A=\varphi_A(A)$
forment un \deleted{\unsure{un|et|vrai}} polygône dans $E$. On \deleted{\unsure{prétend|prétien}} \inserted{dira}
qu'on a un polygône $\Pi$-\uline{épinglé}, ou $\Pi$-\uline{polygône} (le plus souvent, $\Pi$ sera fixé, $E$ et la
% archive sheet number: 50
